\section{本章小结}
\label{sec:opt-ns-summary}

章首问到，尖点处没有唯一斜率时还剩下什么信息。对凸目标，答案是全部全域支撑直线：它们组成次微分，在可微点退化为梯度，在尖点处是一个集合，例如$|x|$在原点的$[-1,1]$、有限最大值在切换点处活跃斜率的凸包。“零属于次微分”与全局最优等价，它是集合层面的判据，不要求算法读到的那一条次梯度为零。失去凸性后，Clarke广义梯度仍能记录局部斜率，但其中含零只是局部极小的必要条件。

只读一条支撑斜率时，负次梯度未必使目标下降，却总能使迭代点靠近最小点。次梯度法的保证因此落在历史最佳点或加权平均点上，目标差为$O(RG/\sqrt T)$，固定步长会留下平台。若非光滑项的邻近映射容易计算，就可以把它完整保留：邻近映射存在唯一、非扩张，其最优条件是在新点读取的次梯度；闭凸集的投影是它的特例，变分不等式和非扩张性都由此得到。邻近梯度法只对光滑项作二次上界近似，在凸情形保持$O(1/T)$的目标差，加速版本达到$O(1/T^2)$，但不再逐轮单调。

没有便宜的邻近映射时，有两条路。光滑化一次性地换成偏差可控的光滑函数，原目标误差分成优化误差与偏差两项，平滑参数同时决定偏差和曲率。代理模型则逐轮重建：MM的上界代理保证原目标单调不增，束方法的下界模型只预测下降，需要查询原目标决定实质步或空步，但它积累的割线给出了停止证书。表~\ref{tab:opt-ns-method-summary}按所需信息、保证与返回对象汇总这些方法。

\begin{table*}[!htbp]
  \centering
  \small
  \begin{tabularx}{\linewidth}{@{}l>{\raggedright\arraybackslash}X>{\raggedright\arraybackslash}X>{\raggedright\arraybackslash}X>{\raggedright\arraybackslash}X@{}}
    \toprule
    方法 & 每轮所需信息与代价 & 主要条件 & 保证的对象 & 返回点 \\
    \midrule
    次梯度法 & 一次函数值与一条次梯度 & 凸，次梯度有界，最小点存在 & 目标差$\leq(R^2+G^2\sum\eta_t^2)/(2\sum\eta_t)$ & 历史最佳点或加权平均点 \\
    邻近点法 & 整个目标的邻近子问题 & 适当、下半连续、凸，最小点存在 & 目标差$O(1/T)$，迭代点收敛 & 最后一点 \\
    邻近梯度法 & 一次梯度与一次邻近映射（回溯时多次） & $f$为$L$光滑，$r$凸且邻近映射易算 & 非凸：最小平方临界残差$O(1/T)$；凸：目标差$O(1/T)$ & 最后一点 \\
    加速邻近梯度法 & 同上，外加外推 & 同上且$f$凸 & 目标差$O(1/T^2)$，不保证单调 & 最后一个记录点 \\
    光滑化 & 替代函数的全部梯度 & 显式的有限最大值结构 & 原目标误差$\leq$优化误差$+\nu\log m$ & 替代问题的近似解 \\
    MM & 构造并近似求解代理 & 合格的全域上界代理 & 原目标单调不增 & 最后一点 \\
    束方法 & 一次查询与一个二次规划 & 有限值凸目标 & 停止证书$e_k+R\lVert\bm s_k\rVert_2$ & 稳定中心 \\
    \bottomrule
  \end{tabularx}
  \caption{本章方法的比较。各行的保证针对不同对象（目标差、临界残差或单调性），复杂度的计数单位也不同（一次查询、一次邻近映射或一个子问题），不能直接按阶数排序。}
  \label{tab:opt-ns-method-summary}
\end{table*}

比较这些方法时要分清三件事。第一，保证的对象：目标差、临界残差、目标单调与参数收敛各不相同，前一项成立不能推出后一项。第二，返回的点：次梯度法的界不针对最后一点，束方法应返回稳定中心而不是候选点。第三，计数单位：一次次梯度查询、一次邻近映射和一个完整子问题的代价相差可能很大，$O(1/\sqrt T)$与$O(1/T)$只有在单位可比时才能比较。

本章的各种保证都假定原问题有最小点；最小点是否存在、是否唯一，仍按第~\ref{chap:optimization}章的下半连续、紧下水平集与强凸条件判断，尖点本身既不帮助也不妨碍这些判断。例如$|x|+\frac\mu2x^2$非光滑却强凸，有唯一最小点；而邻近子问题总有唯一解，并不说明原问题如此。下一章处理约束：指示函数的邻近映射是投影，本章的投影性质在那里成为投影梯度法的基础，但约束的可行性、乘子与对偶证书需要新的工具。第~\ref{chap:opt-stochastic-block}章再处理只能读到带噪声梯度或次梯度的情形，那时本章的确定性误差界要换成期望意义下的界。
